\section{The SQR-SkipResNet-SIREN Backbone: A Deep Dive into the Core Architecture}

This section provides the definitive, mathematically-grounded specifications for each component of the neural network's core structure. Every choice is justified by its contribution to the model's ability to represent complex physical phenomena accurately and to be trained stably.

\subsection{The Sinusoidal Activation (SIREN): The Key to Derivative Fidelity}

The foundational element of the network's layers is the Sinusoidal Representation Network (SIREN) activation function. This choice is not a tunable hyperparameter but a mandatory architectural feature dictated by the physics of the $\mathrm{N}$-body problem. A standard neural network layer performs an affine transformation followed by a non-linear activation. In a SIREN layer, this activation is the sine function$^5$:
\[
\phi_i(\mathrm{x}) = \sin(\mathrm{W}_i \mathrm{x} + \mathrm{b}_i)
\]
where $\mathrm{x}$ is the input to the layer, and $\mathrm{W}_i$ and $\mathrm{b}_i$ are the trainable weight matrix and bias vector, respectively.

The primary motivation for this choice lies in the properties of its derivatives. The derivative of a sinusoidal function is another sinusoid, merely phase-shifted and scaled. For a single neuron with activation $\phi(\mathrm{x}) = \sin(\omega_0 (\mathrm{w} \cdot \mathrm{x} + \mathrm{b}))$, the chain rule gives:
\[
\frac{\mathrm{d} \phi}{\mathrm{d} \mathrm{x}} = \omega_0 \mathrm{w} \cos(\omega_0 (\mathrm{w} \cdot \mathrm{x} + \mathrm{b})) = \omega_0 \mathrm{w} \sin(\omega_0 (\mathrm{w} \cdot \mathrm{x} + \mathrm{b}) + \frac{\pi}{2})
\]

This property is recursive; higher-order derivatives of the network's output with respect to its input remain well-behaved, smooth, sinusoidal functions. This is a stark contrast to common activations like ReLU, whose second derivatives are ill-defined or zero almost everywhere, making them fundamentally unsuited for problems where the loss function depends on second-order derivatives, as is the case for the $\mathrm{N}$-body problem's physics residual.$^1$ The SIREN architecture thus provides an exceptionally strong inductive bias for representing the smooth, continuous solutions and their derivatives required by the governing physical laws. This directly overcomes the spectral bias of standard MLPs, which struggle to represent the high-frequency details inherent in chaotic trajectories and their derivatives.$^2$ By using SIREN, the network is inherently capable of representing the complex natural signals that arise from physical systems.$^8$

\subsection{A Principled Initialization Scheme for SIREN Layers}

The successful training of a deep SIREN is critically dependent on a specific, principled initialization scheme. Naive initialization will cause the distribution of activations to either collapse to zero or explode into saturated regimes, halting the learning process. The confusion between different documents regarding the hyperparameter $\omega_0$ is resolved here by adhering to the original, mathematically-derived SIREN initialization scheme, which is designed to preserve the distribution of activations as they propagate through the network.$^5$

This scheme is not a heuristic but a carefully designed solution to two distinct challenges: representation and trainability. It treats the first layer of the network differently from all subsequent hidden layers.

\begin{enumerate}
    \item
    To optimize network stability during chaotic stages, we dynamically schedule the initial frequency parameter, $\texttt{first\_omega\_0}$, using an exponential decay. This ensures stable representation as complexity increases:

    \begin{lstlisting}[language=Python, caption={Implementation of Dynamic Frequency Scheduling}]
import math

def update_omega0(epoch, total_epochs):
    initial_omega = 30.0   # Initial high-frequency
    final_omega = 10.0     # Stabilizing value
    decay_rate = 5.0       # Controls decay
    return final_omega + (initial_omega - final_omega) * \
           math.exp(-decay_rate * epoch / total_epochs)
    \end{lstlisting}




    \textbf{Rationale}:

    \begin{itemize}
        \item \textbf{Initial high-frequency (30)}: Essential for early stages to overcome spectral bias.
        \item \textbf{Exponential reduction to 10}: Ensures smoother representation in chaotic dynamics, reducing numerical instabilities and promoting stable training.
    \end{itemize}

    \authcomment{Need more detail on this dynamic scheduling for omega\_0.
    Is it based on stages? How do we dynamically adjust it? Based on epochs? Need more details on this.}

    \item \textbf{Subsequent Hidden Layers (for Trainability):} For all other hidden layers, the weights $\mathrm{W}$ are initialized from a
    different uniform distribution, $\mathrm{U}\left[-\frac{\sqrt{6/n_{\text{in}}}}{\text{hidden\_omega\_0}}, \frac{\sqrt{6/n_{\text{in}}}}{\text{hidden\_omega\_0}}\right]$. \\

    The hidden\_omega\_0 parameter is also set to 1. The crucial difference is the scaling of the initialization range by $\omega_0$. This specific scaling is mathematically derived to ensure that the pre-activations (the inputs to the sine function) maintain a standard normal distribution throughout the network's depth. This, in turn, keeps the sine activations in a dynamic, non-saturated range, preventing the frequency of the network's output from growing uncontrollably with depth. This careful, layer-dependent initialization is what makes deep SIRENs trainable by ensuring stable gradient flow.$^1$

    \authcomment{Need confirmation on this hidden layer's detail. Need to make sure the formula is correct according to
    the original SIREN paper?}
\end{enumerate}




\subsection{The Power-Enhanced SkipResNet (SQR-SkipResNet) Formulation}
While the SIREN initialization enables training, pushing to the depths required for chaotic dynamics (e.g., 12 layers) requires
stabilization of the residual connections. This report recommends adopting an advanced ResNet variant with a critical modification:
\textbf{Layer Normalization (LayerNorm)} applied to the skip path before the power term.
This resolves gradient instability while preserving the architecture's benefits for chaotic systems.


The precise mathematical update rules are as follows:

\begin{enumerate}
    \item For odd-numbered layers ($l = 1, 3, 5, \ldots$):
    $$
        \mathbf{x}^{(l)} = f_{\text{sr}}^{(l)}(\mathbf{x}^{(l-1)}) + \left(\text{LayerNorm}(\mathbf{x}^{(l-1)})\right)^p
    $$
    where:
    \begin{itemize}
        \item $f_{\text{sr}}^{(l)}$: SIREN block (SineLayer)
        \item $\text{LayerNorm}(\mathbf{x})$: Normalizes input to zero mean and unit variance per feature
        \item $p = 2$: Power term fixed at 2 (element-wise squaring)
    \end{itemize}

    \item For even-numbered layers ($l = 2, 4, 6, \ldots$):
    $$
        \mathbf{x}^{(l)} = f_{\text{pn}}^{(l)}(\mathbf{x}^{(l-1)})
    $$
\end{enumerate}

\subsubsection*{Critical Rationale for LayerNorm}

A critical enhancement to the standard SQR-SkipResNet is the application of Layer Normalization to the skip path \emph{before} the element-wise squaring operation.
During simulated close encounters, the inputs to these layers (which represent positions and velocities) can vary over orders of magnitude. Without normalization,
the squaring operation would amplify these large values, leading to catastrophic gradient explosions. By normalizing the input on the skip path to have zero mean
and unit variance, this modification stabilizes the training process while preserving the expressive power of the non-linear power term, which is crucial for
modeling the sharp transitions characteristic of close encounters.



\subsubsection*{Power Term Justification}
The power term $p = 2$ remains critical for modeling non-smooth dynamics. For 3-body systems, close encounters produce velocity/acceleration spikes that require quadratic interaction terms to resolve. $\text{LayerNorm}$ ensures this term remains bounded.
\subsection{Finalized Architectural Hyperparameters}

% Table 1: Architectural hyperparameters - can break across pages
\begin{longtable}{|p{3cm}|p{5cm}|p{8cm}|}
\caption{Architectural hyperparameters of the SQR-SkipResNet-SIREN backbone}\label{tab:architectural_hyperparameters}\\
\hline
\textbf{Hyperparameter} & \textbf{Recommended Value} & \textbf{Rationale \& Justification} \\
\hline
\endfirsthead
\caption[]{Architectural hyperparameters of the SQR-SkipResNet-SIREN backbone (continued)}\\
\hline
\textbf{Hyperparameter} & \textbf{Recommended Value} & \textbf{Rationale \& Justification} \\
\hline
\endhead
\hline
\endfoot
\hline
\endlastfoot
Backbone & SQR-SkipResNet with SIREN layers & Superior stability and convergence for non-smooth dynamics, ideal for modeling close encounters in chaotic systems. \\
\hline
Depth & 12 hidden layers (6 SQR-SkipResNet blocks) & Provides a good balance between the expressive power needed for chaotic dynamics and the trainability on modest hardware. \\
\hline
Width & 256 neurons/layer & A strong default that balances model expressivity and computational cost, avoiding the diminishing returns of excessively wide layers.$^1$ \\
\hline
Activation & $\sin(\mathrm{x})$ (SIREN) & The $\mathrm{C}^\infty$-smoothness is mathematically essential for accurately calculating the second-order derivatives required in the physics loss.$^1$ \\
\hline
LayerNorm & elementwise\_affine=True, eps=1e-5 & Stabilizes skip path gradients without restricting feature learning. Applied only to odd layers.\\
\hline
first\_omega\_0 & 30 & A high initial frequency injects a high-frequency basis into the network, helping it overcome spectral bias from the outset.$^1$ \\
\hline
hidden\_omega\_0 & 1 & This value is used within the principled weight initialization formula for hidden layers to ensure a stable distribution of activations, enabling deep network training.$^1$ \\
\hline
Weight Init. (First) & $\mathrm{U}\left[-\frac{1}{n_{\text{in}}}, \frac{1}{n_{\text{in}}}\right]$ & The mathematically-derived initialization for the first SIREN layer to create a random Fourier feature mapping.$^1$ \\
\hline
Weight Init. (Hidden) & $\mathrm{U}\left[-\frac{\text{hidden\_omega\_0}}{6 n_{\text{in}}}, \frac{\text{hidden\_omega\_0}}{6 n_{\text{in}}}\right]$
\authcomment{I believe this formula for initialization is incorrect}
& The principled initialization for all subsequent SIREN layers, designed to preserve activation variance and enable stable training of deep sinusoidal networks.$^1$ \\
\hline
Skip Cadence & Alternating (per SQR SkipResNet) & The alternating pattern of plain and power-enhanced residual layers effectively balances feature transformation and stable gradient flow.$^1$ \\
\hline
Power Term (p) & 2 (element-wise square) & Empirical results demonstrate that $p=2$ provides significant improvements in stability and accuracy without excessive computational overhead.$^1$ \\
\hline
Output Activation & None (Linear) & The final layer must be linear to allow the predicted positions to take on any real value required by the physical system. \\
\hline
\end{longtable}
\FloatBarrier
